% Called after \appendices in main.tex; these are supporting paper sections.
% Keep core definitions, assumptions, and headline results in the main text.
\section{Supplementary derivations and checks}
\label{app:derivations-and-checks}

\subsection{Derivations and theorem assumptions}
\label{app:support-calculations}
% Reserve longer generating-function manipulations, the triangle scalar
% reduction and its admissible domain, and derivative/tangency calculations.
% Refer back to the main equations instead of giving competing definitions.
% Add longer density-derivative calculations and the CM fluctuation-theorem
% checklist: moments, nondegenerate minimum, degree-vector randomness, varying
% parameter sequences, and the effect of divisibility conditioning.
% The main text must still state the assumptions and theoretical status.
\placeholder{[Supporting algebra, density amplitudes, and theorem checks.]}

\subsection{Independent numerical and local checks}
\label{app:numerical-checks}
% Reserve detailed residual/tolerance tables, direct-iteration comparisons,
% local triangle cases, and finite-projection counterexamples where useful.
% Do not label probability-map iteration as a tested finite-graph
% warning-propagation algorithm; that larger experiment is not implemented.
\placeholder{[Supplementary solver checks and small illustrative cases.]}

\section{Experiment settings and sensitivity checks}
\label{app:experiment-details}

\subsection{Configurations and data provenance}
\label{app:configurations}
% Tabulate actual n values, lambda grids, independent repetitions, seeds,
% dataset/stage identifiers, and analysis settings used for the reported run.
% Record tolerances, critical exclusions, event cutoffs, code version, and
% reproducibility commands without reproducing configuration dictionaries.
% Separate existing exploratory runs from allocation screening, calibration,
% validation, and common-lambda paired cohorts. Record vector-level split IDs,
% matching repetitions, positive-count planning, budget caps, and allocation rules.
\placeholder{[Complete experimental settings and reproduction information.]}

\subsection{Additional diagnostics and sensitivity}
\label{app:sensitivity}
% Reserve supplementary degree/erasure panels, outcome plots, and cutoff,
% smoothing, grid, and critical-exclusion checks that support the main claims.
% Give each retained diagnostic its methodological reference and stated purpose.
% Omit unused exploratory exports instead of treating all generated files as
% necessary report figures.
% Add bootstrap/cluster details, occurrence-fit checks, goodness-of-fit
% calibration, full cutoff panels, sparse-tail counts, and secondary ECDFs.
\placeholder{[Distributional checks, uncertainty, and cutoff sensitivity.]}
